\section{Conclusion}
\label{sec:conclusion}

\ALUCLU is an attempt to make memory trade-offs architectural objects rather
than footnotes. Dense recurrence, exact recency, sparse exceptional records,
and compressed episodes are assigned separate bounded states, update rules,
and failure modes. Learned fusion lets the model combine those states at each
token, while the public recurrent interface keeps state ownership explicit.

The reference implementation establishes a precise baseline for the idea. It
defines inclusive write-then-read causality, deterministic cache replacement,
an auditable episodic life cycle, small-core factor compression, a
compression-error ledger, conditional mass conservation, and concrete
persistent-byte formulas. The automated tests and small CPU measurements show
that these mechanisms are executable and internally consistent within the
tested domain.

The work is intentionally unfinished at the level where strong empirical
claims begin. The central research questions now have measurable forms:
whether residual surprise selects useful long-range evidence, whether the
episodic lane improves quality per state byte, whether token-wise fusion
specializes by dependency type, and whether faithful fused kernels move the
quality-throughput-memory Pareto frontier. The evaluation program in
\Cref{sec:evaluation} is designed to answer those questions without confusing
mechanical correctness with model quality or reference code with optimized
systems.

The practical thesis is simple:
\begin{quote}
\centering
\emph{Layer the context, bound the memory, expose the loss, and measure the
result.}
\end{quote}
If later experiments reject a lane or a routing rule, the architecture's
modularity makes that a useful result rather than an untraceable failure. If
they validate the composition, the same specification provides the semantic
oracle required to engineer it aggressively without losing correctness.
